% ======================================================================
%  1. Introduction
% ======================================================================
\section{Introduction}
\label{sec:intro}

\subsection{Scientific and laboratory context}
\label{sec:intro-context}

The sciNeurotech Lab studies how targeted electrical stimulation of the nervous
system can be used to drive and restore motor function. In a typical experiment,
a stimulation configuration~$x$, an electrode, a current amplitude, a
frequency, a pulse shape is applied, and a biological response $f(x)$ is
recorded, for instance an evoked muscle activation or a cortical readout. The
map $x\mapsto f(x)$ is \emph{unknown}, \emph{noisy}, and above all
\emph{expensive}: every evaluation consumes experimental time on a living
subject, so the number of stimulations that can be tried is small.

These are exactly the conditions under which Bayesian optimization (BO) is the
tool of choice. A Gaussian process (GP) surrogate models the response from the
handful of stimulations tried so far, and an acquisition function decides which
configuration to try next, balancing exploration of uncertain regions against
exploitation of promising ones. GP-based BO has already been used successfully
for neurostimulation in the group and elsewhere; this internship takes the
method as given and asks a more structural question about \emph{how} it should be
set up.

\subsection{One phenomenon, several decision problems}
\label{sec:intro-onephenomenon}

The starting observation is that a single biological phenomenon $f$ is rarely of
interest for one fixed goal. The very same response surface can be interrogated
in several ways, depending on the scientific question:
\begin{itemize}
  \item \textbf{Target matching:} find a stimulation that produces a
        \emph{prescribed} response level $T$, i.e.\ solve
        $\min_x \abs{f(x)-T}$.
  \item \textbf{Multi-objective optimization:} arbitrate between several
        competing responses $F(x)=(f_1(x),\dots,f_m(x))$ for example a hand,
        an elbow and a biceps activation for which there is no single
        optimum but a trade-off (Pareto) surface.
  \item \textbf{Constrained optimization:} maximize one response subject to
        safety or physiological constraints $c(x)\le 0$.
  \item \textbf{Prior-guided optimization:} exploit prior knowledge that
        effective stimulations are likely to lie in a particular region of the
        parameter space.
\end{itemize}
The phenomenon $f$ does not change from one task to the next; what changes is the
\emph{decision problem} built on top of it. This is the thread that runs through
the whole report.

\subsection{Separation of concerns}
\label{sec:intro-soc}

Each downstream task can be described by a \emph{task transformation} $\prm$
applied to the phenomenon. Target matching uses $\prm(z)=-\abs{z-T}$;
scalarized multi-objective optimization uses a scalarization
$\prm_\lambda:\R^m\to\R$; prior-guided optimization reweights the search by a
spatial prior $\pi(x)$. There are then two fundamentally different places where
this task information can be injected (Figure~\ref{fig:soc-diagram}):
\begin{description}
  \item[Separated formulation (\RIGHT).] The GP models the raw phenomenon $f$,
        and the task transformation $\prm$ enters only the acquisition /
        decision rule.
  \item[Task-in-model formulation (\WRONG).] The GP is fitted directly to the
        transformed quantity $\prm\circ f$, so the task is baked into the object
        the model must learn.
\end{description}

\begin{takeaway}
The phenomenon may stay the same while the scientific task changes. The
central principle studied in this report is: \emph{model the phenomenon $f$, and
encode the task in the acquisition whenever possible.}
\end{takeaway}

The names \RIGHT{} and \WRONG{} are used throughout as convenient labels for the
two architectures; they should not be read as a blanket claim that non-trivial
GP means or transformed models are always harmful. As we make precise in
Section~\ref{sec:contrib-prior}, a genuine \emph{functional} prior over the whole
response surface can be theoretically beneficial. The claim of this work is
narrower and structural: task information of the kind that arises in the problems
above - a target level, a preference weight, a probable optimizer location -
can generally be moved to the decision level, and doing so preserves assumptions
that folding it into the model can destroy.

The main research question is therefore:
\begin{center}
\emph{When does putting task information into the statistical model change the
assumptions on which Bayesian optimization relies?}
\end{center}

\subsection{Research questions}
\label{sec:intro-questions}

The work is organized around three sub-questions.
\begin{enumerate}
  \item \textbf{Task-induced model misspecification.} Does encoding the task in
        the model take the optimized function outside the function class assumed
        by the kernel, and if so, how severely?
  \item \textbf{Preservation of confidence and regret guarantees.} Does the
        separated formulation retain the confidence bounds and regret guarantees
        of standard GP-UCB, and through what mechanism?
  \item \textbf{Generalization beyond target matching.} Does the same
        separation principle, and the same pushforward-confidence reasoning,
        extend to multi-objective and prior-guided optimization?
\end{enumerate}

\subsection{Contributions}
\label{sec:intro-contributions}

The provenance of every result is made explicit throughout, using the following
colour code, which also appears next to each theorem statement:

\begin{center}
\begin{tabular}{@{}ll@{}}
  \legendchip{colLit}{established literature} & standard, imported without change \\[3pt]
  \legendchip{colSpec}{literature, specialized here} & a known result re-derived / specialized for our setting \\[3pt]
  \legendchip{colContrib}{personal contribution} & a personal theoretical contribution \\[3pt]
  \legendchip{colExp}{personal experiment} & a personal experimental contribution \\[3pt]
  \legendchip{colSynth}{interpretation / synthesis} & an interpretation or synthesis proposed here \\[3pt]
  \legendchip{colOpen}{open question} & an open or conjectural statement \\
\end{tabular}
\end{center}

\noindent My personal contributions are the following.
\begin{itemize}
  \item \textbf{A separation-of-concerns reading of GPBO} that unifies target
        matching, multi-objective optimization and prior-guided optimization
        under a single question - where task information should enter -
        together with a \RIGHT/\WRONG{} taxonomy
        (Sections~\ref{sec:contrib-tm}--\ref{sec:contrib-prior}).
        \hfill\legendchip{colSynth}{interpretation / synthesis}
  \item \textbf{A quantitative approximation barrier} for the target-matching
        objective $g=\abs{f-T}$: I show that a transverse target crossing
        creates a cusp, that under smooth-kernel assumptions this places $g$
        outside the RKHS, and - the stronger statement - that the best
        approximation of $g$ by an RKHS ball of radius $B$ cannot decay faster
        than a power of $B$ (Section~\ref{sec:contrib-tm-barrier},
        Appendix~\ref{app:tm}).
        \hfill\legendchip{colContrib}{personal contribution}
  \item \textbf{Positive separation guarantees.} I show that retaining the GP on
        $f$ and pushing its confidence interval through the task yields a
        GP-UCB-type regret bound for target matching, and I extend this
        \emph{pushforward-confidence} argument to Lipschitz scalarizations in the
        multi-objective setting (Sections~\ref{sec:contrib-tm-pos},
        \ref{sec:contrib-mo}, Appendices~\ref{app:tm},~\ref{app:mo}).
        \hfill\legendchip{colContrib}{personal contribution}
  \item \textbf{A separation-of-concerns treatment of functional priors}, introducing a prior-weighted $\alpha$-UCB acquisition with convergence guarantees, and identifying two failure modes when the prior is embedded in the surrogate: persistent global misspecification despite good local fit, and self-confirming feedback under adaptive sampling. (Section~\ref{sec:contrib-prior},
        Appendix~\ref{app:prior}).
        \hfill\legendchip{colContrib}{personal contribution}
  \item \textbf{A separation-of-concerns interpretation of random-scalarization
        multi-objective BO,} making explicit why a single vector observation
        $F(x_t)$ informs all future preferences under a shared objective-wise
        model, whereas independent scalarized surrogates cannot share this
        information (Section~\ref{sec:contrib-mo-sharing}).
        \hfill\legendchip{colSynth}{interpretation / synthesis}
  \item \textbf{Three controlled experimental studies} (target matching,
        multi-objective, prior-guided) that operationalize the \RIGHT/\WRONG{}
        contrast on standard benchmarks with matched budgets, shared seeds and
        paired tests (Sections~\ref{sec:contrib-tm-exp},
        \ref{sec:contrib-mo-exp}, \ref{sec:contrib-prior-exp},
        Appendix~\ref{app:exp}).
        \hfill\legendchip{colExp}{personal experiment}
\end{itemize}

\noindent I am careful \emph{not} to overclaim. In particular, the approximation
barrier is not by itself a lower regret bound for task-in-model GP-UCB, and I do
not claim that the task-in-model method necessarily fails or does not converge;
these points are stated as open questions
(Section~\ref{sec:discussion}).

\subsection{Organization of the report}
\label{sec:intro-organization}

Section~\ref{sec:background} gives a selective, argumentative state of the art,
covering only the material needed for the contributions. Section~\ref{sec:contrib-tm}
develops target matching, from the qualitative regularity obstruction to the
quantitative barrier and the positive separation guarantee, followed by
Experiment~A. Section~\ref{sec:contrib-mo} generalizes the principle to
multi-objective optimization (pushforward confidence, information sharing across
preferences), followed by Experiment~C. Section~\ref{sec:contrib-prior} treats
prior-guided optimization (spatial versus functional priors, self-confirmation),
followed by Experiment~B. Section~\ref{sec:discussion} discusses limitations and
open problems, and Section~\ref{sec:conclusion} concludes. Appendices~%
\ref{app:background}--\ref{app:exp} collect the mathematical background, the full
proofs, and the complete experimental configurations.

% --- SoC diagram --------------------------------------------------------
\figauto[0.9\linewidth]{figures/Theory/soc_diagram.png}{%
  \textbf{Separation of concerns.} Left (\RIGHT): the GP models the raw
  phenomenon $f$; the task transformation $\prm$ acts only at the acquisition
  level. Right (\WRONG): the transformation $\prm$ is baked into the object the
  GP must learn, $\prm\circ f$. This contrast reappears in target matching,
  multi-objective and prior-guided optimization. }{fig:soc-diagram}
